\documentclass[11pt,a4paper]{article}

\usepackage[utf8]{inputenc}

\usepackage[T1]{fontenc}

\usepackage[polish]{babel}

\usepackage{lmodern}

\usepackage{geometry}

\geometry{a4paper,margin=27mm}

\usepackage{microtype}

\usepackage[hidelinks]{hyperref}

% Portfolio: MSC-AR-P23; source: prospective ARC-XR research scope.

\title{Adaptacja poziomu wyzwania XR a dynamika subiektywnego doświadczenia flow}

\author{Mateusz Pomianek}

\date{Wersja 0.3.0, 21 listopada 2026}

\begin{document}

\maketitle

\begin{abstract}

Przedmiotem opracowania jest Adaptacja poziomu wyzwania XR a dynamika subiektywnego doświadczenia flow. Czy adaptacja trudności jest związana ze zmianą doświadczanego flow w trakcie zadania? Porównane zostaną stała i adaptowana trudność w zadaniu XR, przy wspólnym pomiarze poprawności oraz doświadczenia flow. Głównym punktem końcowym jest wynik narzędzia mierzącego flow po zadaniu oraz poprawność zachowania. Opis ma charakter wstępny i nie przedstawia niezweryfikowanych wyników.

\end{abstract}

\section{Wprowadzenie}

Rozwój adaptacyjnych systemów XR stwarza możliwość badania interakcji człowieka z informacją przestrzenną, ale samo zastosowanie AR, VR lub modeli sztucznej inteligencji nie przesądza o przewadze poznawczej czy ergonomicznej. W rozpatrywanym zagadnieniu, Adaptacja poziomu wyzwania XR a dynamika subiektywnego doświadczenia flow, kluczowe jest rozdzielenie właściwości systemu od zmienności zachowania użytkownika. Porównanie należy zaprojektować tak, aby mierzyć jednoznaczne zjawisko, a nie tylko ogólną ocenę technologii.

\section{Cel pracy i pytanie badawcze}

Czy adaptacja trudności jest związana ze zmianą doświadczanego flow w trakcie zadania? Projekt wymaga krytycznego przeglądu aktualnego stanu badań oraz zdefiniowania nowości w stosunku do innych tematów portfolio XR.

\section{Zakres i odrębny wkład naukowy}
Subiektywny konstrukt flow w zadaniach XR, niezależny od poprawności transferu umiejętności badanej w P08. Należy użyć uprzednio określonego, zwalidowanego narzędzia. Przyczynowy kontrast flow jako odrębnego doświadczenia, a nie ponowne opisanie efektu poprawności P08.
\section{Hipotezy badawcze}
\paragraph{H1.} Przy wyrównanej ekspozycji adaptacja trudności zmienia średni wynik flow względem stałego poziomu wyzwania.
\paragraph{H2.} Efekt flow nie jest całkowicie wyjaśniony poprawnością wykonania zadania lub poziomem trudności początkowej.
\paragraph{H3.} Zmiana flow wykazuje inną dynamikę w kolejnych blokach niż zmiana trafności wykonania, uzasadniając osobny konstrukt.
\section{Projekt badania i plan eksperymentów}
\subsection{A. Operacjonalizacja i kwalifikacja danych}
Subiektywny konstrukt flow w zadaniach XR, niezależny od poprawności transferu umiejętności badanej w P08. Należy użyć uprzednio określonego, zwalidowanego narzędzia. Jednostka wnioskowania: uczestnik × sesja z oceną flow po zadaniu. Przed badaniem należy określić źródła, wersje urządzeń, czas akwizycji, etykiety i niezależną referencję. Nie liczyć zależnych okien i klatek jako niezależnych prób.
\subsection{B. Eksperyment porównawczy}
Randomizować warunki trudności i równoważyć zadania; zebrać flow po całych blokach zamiast zakłócać immersję ciągłym sondowaniem. Wynik pierwszorzędowy: różnica wyniku wcześniej określonej skali flow: adaptowana vs stała trudność przy wyrównanym czasie. Zamrozić losowanie warunków, progi, mianownik punktu końcowego i kryteria wyłączeń bez dostępu do oceny końcowej.
\subsection{C. Kontrole i analiza odporności}
Kontrolować wyniki zadania, frustrację i zmęczenie; nie interpretować flow jako stanu niskiego obciążenia. Nie publikować drugiej pracy na tej samej próbie, jeśli analiza jest jedynie wtórną tabelą P08 bez osobnej hipotezy i raportu nakładania. Utrata sygnału, abstencja, przerwanie zadania i dane brakujące są odrębnymi kategoriami wyniku.
\subsection{D. Walidacja na niezależnym materiale}
Walidacja musi zachować podział po uczestniku, sesji, scenie lub zadaniu. Model mieszany wyników zwalidowanej skali i bloku zadania; oddzielna analiza wyników behawioralnych. Raportować wyniki negatywne i rozbieżności z założeniami.
\section{Estymand, analiza statystyczna i decyzje}
różnica wyniku wcześniej określonej skali flow: adaptowana vs stała trudność przy wyrównanym czasie. Jednostka: uczestnik × sesja z oceną flow po zadaniu. Model mieszany wyników zwalidowanej skali i bloku zadania; oddzielna analiza wyników behawioralnych. H1 jest planowanym kontrastem głównym; H2 i H3 pozostają pomocnicze do zatwierdzenia hierarchii testów oraz kontroli wielokrotności. Wartości marginesów, mocy, minimalnego efektu istotnego praktycznie i liczebności mają wynikać z pilotażu, a nie z post hoc wyboru wyniku.
\section{Stan badań i granica nowości}
\begin{itemize}
\item XR personalization methods exist; a claim needs a defined multi-session or transfer contrast. \cite{xrprior1}.
\item Physiology-driven VR adaptation has been demonstrated; independent behavioural gains require direct controlled trials. \cite{xrprior2}.
\end{itemize}
Przyczynowy kontrast flow jako odrębnego doświadczenia, a nie ponowne opisanie efektu poprawności P08. Odrębność na tle opublikowanych metod i pozostałych artykułów wymaga osobnego audytu.
\section{Warunki bezpieczeństwa, sprzętu i zgodności}
To prospektywna koncepcja, nie zgoda na eksperyment, prerejestracja lub wynik empiryczny. Praca z uczestnikami, interwencje stresowe i pomiar biometryczny wymagają przeglądu etycznego, świadomej zgody i ochrony danych. Quest 3 w wersji standardowej nie posiada eye trackingu, a Tobii Eye Tracker 5 jest urządzeniem ekranowym; nie zastępować jednego stanowiska drugim.
\section{Materiał i metody}

Porównane zostaną stała i adaptowana trudność w zadaniu XR, przy wspólnym pomiarze poprawności oraz doświadczenia flow. Wymagane jest zdefiniowanie jednostki analizy, zasad doboru przypadków, wersji scenariusza i źródeł pomiaru. Rozkład warunków oraz reguły wykluczenia prób należy określić przed pomiarem. W badaniach z udziałem ludzi niezbędna jest uprzednia właściwa ocena etyczna i świadoma zgoda. W testach wyłącznie symulacyjnych należy oddzielić prawdę symulacyjną od obserwacji empirycznej.

\subsection{Wynik główny i kryteria jakości}

Wynik narzędzia mierzącego flow po zadaniu oraz poprawność zachowania. Pomiar powinien zachowywać jawne jednostki, znaczniki czasu i reguły oceny poprawności. Przy modelach predykcyjnych stosować rozłączne partycje oceny; przy eksperymentach interwencyjnych oddzielić skuteczność predykcji od przyczynowego wpływu interfejsu.

\section{Plan analizy}

Zestawić zależności nieliniowe między trudnością, biegłością i flow, analizując zmiany w czasie. W wynikach należy raportować wielkość efektu lub błąd estymacji z przedziałem niepewności, liczbą obserwacji i odsetkiem odrzuconych prób. Analizy dodatkowe wymagają jasnego oznaczenia jako eksploracyjne. Nie wybierać progów decyzyjnych ani modelu na podstawie obserwacji ze zbioru testowego.

\section{Wyniki}

W obecnej wersji nie raportuje się wyników pomiarów, tabel statystycznych ani parametrów wyznaczonych z nieudokumentowanych danych. To kierunek badawczy, a nie zatwierdzony protokół badania ani zgłoszony artykuł.

\section{Dyskusja}

Interpretacja przyszłych wyników musi uwzględnić alternatywne mechanizmy i zależność wniosków od jakości danych, warunków zadania oraz ograniczeń sprzętu. Wartość potencjalnego rozwiązania zależy nie tylko od trafności modelu, lecz także od jego opóźnienia, stabilności i wpływu na zachowanie człowieka. Wyniki testów syntetycznych nie mogą być automatycznie uogólniane na populację uczestników.

\section{Ograniczenia badania}

Flow nie może być utożsamiony z wysoką wydajnością ani ze stanem niskiego obciążenia poznawczego. Konstrukt flow, nie skuteczność wykonania; drugi artykuł tylko przy osobnej hipotezie i niezależnej analizie. Żadna część szkicu nie dowodzi przeprowadzenia zbierania danych, uzyskania zgód etycznych ani gotowości manuskryptu do zgłoszenia.

\section{Wnioski}

Przedstawiono problem, kontrast i sposób jego planowanej oceny. Rozstrzygnięcie postawionego pytania wymaga przygotowania protokołu, zweryfikowanych źródeł oraz poprawnie udokumentowanych wyników. Bez tych materiałów nie można sformułować konkluzji empirycznych.

\begin{thebibliography}{9}
\bibitem{xrprior1} Adapting to the User: A Systematic Review of Personalized Interaction in VR (2025), arXiv:2510.13123. \url{https://arxiv.org/abs/2510.13123}.
\bibitem{xrprior2} Towards Intelligent VR Training: A Physiological Adaptation Framework for Cognitive Load and Stress Detection (2025), UMAP 2025. \url{https://doi.org/10.1145/3699682.3727575}.
\end{thebibliography}
\end{document}

